\documentclass[conference]{IEEEtran}
\IEEEoverridecommandlockouts
% The preceding line is only needed to identify funding in the first footnote.
\usepackage{cite}
\usepackage{amsmath,amssymb,amsfonts}
\usepackage{algorithmic}
\usepackage{graphicx}
\usepackage{textcomp}
\usepackage[dvipsnames]{xcolor}
\usepackage{booktabs}
\usepackage{multirow}
\usepackage{colortbl}
\usepackage[normalem]{ulem}
% Highlighting helpers for the quantitative-comparison table:
% best (green, bold), second (light green, underlined), third (yellow).
\newcommand{\best}[1]{\cellcolor{ForestGreen!20}\textbf{#1}}
\newcommand{\second}[1]{\cellcolor{LimeGreen!20}\uline{#1}}
\newcommand{\third}[1]{\cellcolor{yellow!20}#1}
\def\BibTeX{{\rm B\kern-.05em{\sc i\kern-.025em b}\kern-.08em
    T\kern-.1667em\lower.7ex\hbox{E}\kern-.125emX}}

\begin{document}

\title{SHARC-GS: Scale-aware Hierarchical Anchor-based Reproducible
Compression for Gaussian Splatting}



\maketitle

\begin{abstract}
3D Gaussian Splatting (3DGS) achieves high-quality real-time novel-view
rendering through explicit Gaussian primitives, but the scale of
primitive attributes makes model size and storage/transmission overheads
grow rapidly in large-scale UAV aerial scenes, becoming a key bottleneck
in practical deployment. This contribution proposes SHARC-GS (Scale-aware
Hierarchical Anchor-based Reproducible Compression for Gaussian
Splatting) on top of the HAC++ Anchor--Hash compression framework, with
two complementary modifications. First, a scale-aware hierarchical
context augments fine-grained hash features with parent-grid features
and a three-level observation-distance label, so the entropy model can
distinguish anchors at near, mid, and far scales. Second,
decoder-reproducible content-adaptive quantization recomputes each
quantization step on the fly from structural quantities that the decoder
can also reconstruct, retaining the rate--distortion benefits of
content-adaptive quantization without any additional per-anchor step
data. On three large-scale UAV scenes, the default $\lambda=0.004$
operating point requires 17.151/15.488/6.355 MB at PSNR
29.296/28.629/28.563 dB; compared with the HAC++ reference points, it
improves PSNR by 0.25 dB on 4-28 and reaches its best PSNR on 1-100,
while reducing size by 11.4\% on 3-07. These per-scene outcomes are
scene dependent and do not establish a uniform rate--distortion
advantage; we therefore report the full matched-budget curve and
Bjontegaard metrics rather than a single operating point. The value of
the codec lies in the combination of a self-contained bitstream with
zero symbol/CDF mismatch, decoder reproducibility, and runtime
efficiency.
\textbf{[TODO: B1 BD-Rate summary -- PSNR/SSIM/LPIPS BD-Rate \& BD-PSNR on
matched budget]}
All bitstreams pass standalone decoding verification in a fresh process:
symbol, quantization-step, and integer-CDF mismatch counts are all zero.
\end{abstract}

\begin{IEEEkeywords}
Gaussian Splatting, compression, entropy coding, content-adaptive
quantization, scale-aware context, decoder reproducibility
\end{IEEEkeywords}

\section{Introduction}
3DGS represents a scene with a large number of anisotropic Gaussian
primitives and renders them through a visibility-aware pipeline,
which supports high-fidelity novel-view synthesis at real-time speed. However,
geometry, covariance, opacity, and appearance attributes are stored
explicitly per primitive, so model size grows rapidly with scene scale.
UAV aerial scenes are particularly demanding: camera heights and
observation distances span a wide range and flight paths overlap
heavily, so a model must preserve both large-scale structure and fine
local detail.

In UAV mapping, urban modeling, and facility inspection, 3D models are
often transmitted over constrained bandwidth or resident on edge devices
for long periods. Beyond compression ratio, bitstream self-containment,
decoding reproducibility, and inference efficiency therefore decide
whether a solution can actually be deployed. This contribution is
designed around these goals.

Anchor-based structured representations such as Scaffold-GS
\cite{lu2024} generate view-adaptive Gaussians from anchors and
substantially reduce primitive redundancy. HAC \cite{chen2024} and HAC++
\cite{chen2025} further exploit the mutual information between anchors
and a structured hash grid, predicting anchor-attribute distributions
with a conditional entropy model and entropy coding, and achieve
orders-of-magnitude size reduction. HAC++ is the reference framework of
this contribution.

We propose SHARC-GS by enhancing the HAC++ pipeline on two levels.
First, we upgrade the single-scale spatial context to a scale-aware
hierarchical context, so the entropy model explicitly exploits
cross-scale structural priors and observation-distance information.
Second, we realize content-adaptive quantization as a
decoder-reproducible formula, so quantization steps adapt to local
content without any per-anchor side information. Both modules reuse the
HAC++ main pipeline without modifying anchor generation, Gaussian
expansion, or rendering.

Our contributions are as follows: (1) a scale-aware hierarchical
Anchor--Hash context that jointly conditions the entropy model on
fine-grained hash features, parent-grid features, and a three-level
observation-distance label; and (2) decoder-reproducible
content-adaptive formula quantization, where quantization steps are
generated on the fly from reconstructable structural quantities and the
bitstream carries no per-anchor step array. Both modules are realized as
a self-contained bitstream whose symbols, quantization steps, and
integer CDFs are verified by standalone decoding (Section~\ref{sec:bitstream}).

Training, encoding, and decoding share the same context-construction
rule and quantization formula; the encoder produces a self-contained
bitstream, and the decoder recovers identical context, quantization
steps, and entropy parameters from the bitstream alone in a fresh
process.

\section{Related Work}
Since the introduction of 3DGS \cite{kerbl2023}, compression has become
central to its large-scale adoption. Existing work falls into three
broad lines. Raw-Gaussian compression prunes, quantizes, or
entropy-codes the Gaussian primitives directly, with limited size
reduction. Anchor-based structured methods such as Scaffold-GS
\cite{lu2024} generate Gaussians implicitly from anchors and inherently
remove inter-primitive redundancy. Context-modeling methods such as HAC
\cite{chen2024} and HAC++ \cite{chen2025} predict anchor-attribute
distributions from hash-grid context and stand among the strongest
context-based compression methods.

In a compression pipeline, quantized symbols are mapped by entropy
coding to code words close to the information-theoretic limit, and the
actual code length depends on how well the probability model matches the
true distribution; conditional probability estimation therefore
determines the compression ceiling. The mask-aware factorization of
HAC++ decomposes the anchor joint distribution by attribute and offset
factors and sharpens estimation with hash-grid context, providing a
clear extension point for further improvement.

Closely related are adaptive quantization and content-aware coding.
HAC++ introduces a learnable adaptive quantization module for anchor
attributes to improve fidelity in high-precision regions, while
contextgs and PC-GS improve attribute coding from context modeling and
progressive compression perspectives, respectively. Together these works
show that sharper conditional modeling and finer bit allocation are the
two main lines toward better rate--distortion performance.

The methods compared in Table~\ref{tab:cmp} cover context modeling,
structure pruning, vector quantization, and progressive coding.
Anchor-based methods (Scaffold-GS, 3DGS, gsplat) keep hundreds of MB to
preserve the full structure, while methods that trade quality
aggressively for small size (e.g., RDO-gaussian) sacrifice substantial
quality. Across these lines, no method couples scale-aware entropy
modeling with decoder-reproducible content-adaptive quantization inside a
self-contained bitstream, and none evaluates multi-view or temporal
consistency on large-scale aerial scenes. This is the gap we address.

\section{Method Description}
Our method builds on the HAC++ Anchor--Hash compression framework. For
the $i$-th anchor, $x^{(i)}$ denotes its position, $f^{(i)}$ its
feature, $s^{(i)}$ its scaling, and $o^{(i,k)}$, $m^{(i,k)}$ denote
$K$ candidate offsets and their learned masks; $\theta^{(i)}$ collects
the encoded attributes.

We pose the codec as a constrained rate--distortion problem. Given a
target bitstream budget $B$ and a clean-decoder contract, the encoder
minimizes decoding distortion subject to (i) the bitstream being
decodable by a fresh process that never sees training metadata, and (ii)
quantization being reproducible on the decoder side. $\lambda$ is the
Lagrange multiplier of the rate--distortion objective, and $N_\theta$
normalizes the rate term so that different anchor counts are comparable.
Within this setup, I1 sharpens how the probability model conditions on
spatial scale, and I2 turns content-adaptive quantization into a
decoder-reproducible rule. We describe the two modules below.

The entropy model factorizes the attribute joint distribution in a
mask-aware way, so offset dimensions turned off by the mask do not
participate in coding; the probability of a quantized symbol is obtained
by integrating the density over its quantization cell. A differentiable
rate proxy sums the weighted log-probabilities of symbols and adds the
hash-parameter term. The base quantization step follows a tanh-bounded
form, as in
\begin{equation}
Q^{(i,t)} = Q_{0}^{(i,t)} \cdot \bigl(1 + \tanh\bigl(r^{(t)}\bigr)\bigr),
\label{eq:basequant}
\end{equation}
\noindent where $Q_0^{(i,t)}$ is the per-anchor base step for field $t$
and $r^{(t)}$ is a field-wise learned offset that HAC++ stores with the
model. In Eq.~\eqref{eq:qstep} we keep this tanh family but replace the
stored offset $r^{(t)}$ with per-anchor, per-field logits
$z^{(i,t)}$, so the two equations are the same quantization-step family
under different conditioning: HAC++ transmits the offset, whereas
SHARC-GS recomputes it from the decoder-reproducible descriptor
$u^{(i)}$. The coefficient $\alpha$ scales the modulation; we use
$\alpha=0.35$ throughout.
and the conditional joint distribution is factorized as
\begin{equation}
\begin{split}
p(\theta|C) = \prod & p_{f}\bigl(f^{(i)}|c^{(i)}\bigr) \cdot
p_{s}\bigl(s^{(i)}|c^{(i)}\bigr) \\
& \cdot \prod
\bigl[p_{o}\bigl(o^{(i,k)}|c^{(i)}\bigr)\bigr]^{m^{(i,k)}}.
\end{split}
\label{eq:entropy}
\end{equation}
The training objective combines rendering loss, scaling-volume
regularization, and a rate--distortion term
\begin{equation}
\mathcal{L} = (1-\beta)\,L_{1} + \beta\,(1-\text{SSIM}) +
\eta\,\frac{1}{N}\sum_{i}\prod_{d}s(i,d) + \lambda\,\frac{R}{N_{\theta}}.
\label{eq:loss}
\end{equation}
The default values are $\beta=0.2$ and $\eta=0.01$, while $\lambda$
controls the rate--distortion trade-off. Training, encoding, and
decoding share the same conditional probability model and quantization
path, so the two ends reconstruct exactly the same entropy parameters.

\begin{figure*}[t]
\centering
\includegraphics[width=0.92\textwidth]{fig1.png}
\caption{Overview of SHARC-GS: the encoder builds a self-contained
bitstream with hierarchical context and reproducible formula
quantization; the decoder reconstructs identical state in a fresh
process. The bottom strip is the three-scene standalone decoding
verification.}
\label{fig:overview}
\end{figure*}

\begin{table}[t]
\caption{Notation}
\centering
\begin{tabular}{@{}ll@{}}
\toprule
\textbf{Symbol} & \textbf{Meaning} \\
\midrule
\(x^{(i)},f^{(i)},s^{(i)}\) & position, feature, and scaling of the $i$-th anchor \\
\(o^{(i,k)},m^{(i,k)}\) & $k$-th candidate offset and its mask \\
\(h(x)\) & interpolated hash-grid feature at $x$ \\
\(\bar{x}^{(i)}\) & parent-grid aligned coordinate of the anchor \\
\(\ell^{(i)}\) & three-level observation-distance label (near/mid/far) \\
\(c^{(i)}\) & conditioning vector of the entropy model \\
\(Q^{(i,t)},Q_{0}^{(i,t)}\) & quantization step and base step of field $t$ \\
\(\alpha\) & quantization-step adjustment magnitude (0.35) \\
\(\lambda\) & rate--distortion weight \\
\(R,N_{\theta}\) & rate proxy and number of attribute parameters \\
\bottomrule
\end{tabular}
\label{tab:notation}
\end{table}

\subsection{Scale-Aware Hierarchical Anchor--Hash Context}
The base condition follows HAC++: the interpolated fine-grid hash
feature $h(x^{(i)})$ at the anchor position retains local structure. To
add a coarse-scale prior, we align the anchor to a parent grid whose
spacing is four times the fine-grid spacing and query the same hash
table for the parent feature $h(\bar{x}^{(i)})$. Because the parent grid
reuses the already-trained hash table, this coarse-scale prior adds no
parameters and no bitstream payload, whereas adding a new hash
resolution level would enlarge the hash parameters. We also split
anchors into near, mid, and far levels by the average observation
distance accumulated during training, with cutoffs at the 33\% and 66\%
percentiles of all valid observation distances, yielding the three-level
distance label $\ell^{(i)}$. Three levels keep the label to 2 bits per
anchor, and the cutoffs adapt to the data. The final context is
constructed identically for training, encoding, and decoding:
\begin{equation}
c^{(i)} = \bigl[\, h\bigl(x^{(i)}\bigr),\; h\bigl(\bar{x}^{(i)}\bigr),\;
\text{onehot}\bigl(\ell^{(i)}\bigr) \,\bigr].
\label{eq:context}
\end{equation}
The value of scale information lies in aligning conditional modeling
with scene structure: nearby anchors are detail-rich and more sensitive
to probability estimation, while distant anchors are structurally sparse
and tolerate coarser modeling; the parent feature adds large-range
context beyond a single anchor neighborhood. Component ablations are
consistent with its independent contribution (Section~\ref{sec:ablation}).

The hierarchical context is enabled only after the 12,000-th training
iteration to avoid conditioning on unstable early geometry; it changes
only the conditional input of the probability model, not Gaussian
generation or rendering. In the bitstream contract, each entropy-coded
anchor needs only 2 packed bits for its distance label: with
520,771/519,606/251,351 anchors on the three scenes, the label payload
is about 0.06--0.13 MB per scene and is negligible.

The header records the version, field ranges, anchor count, and a
SHA-256 checksum, from which the decoder recovers exactly the same
conditioning. The distance label is therefore part of the formal
bitstream rather than an implicit training-time dependency.

\subsection{Decoder-Reproducible Content-Adaptive Formula Quantization}
Uniform quantization uses one step for every region and cannot serve
flat and detailed regions equally well: flat regions can be quantized
more coarsely to save bits, while detailed regions need finer
quantization to preserve quality. Content-adaptive quantization adjusts
the step according to regional complexity and is an effective way to
improve rate--distortion performance.

The key design is decoder reproducibility: the step is not decided by
the encoder and transmitted, but regenerated by the decoder from
structural quantities it can also reconstruct. For each anchor we
extract four structural quantities---local anchor density, predicted
scaling anisotropy, predicted offset energy, and active mask ratio---and
feed them into a small complexity network that outputs adjustment logits
$z^{(i,t)}$ for the Feature, Scaling, and Offset fields:
\begin{equation}
Q^{(i,t)} = Q_{0}^{(i,t)} \cdot \bigl(1 + \alpha\tanh\bigl(z^{(i,t)}\bigr)\bigr),
\qquad \alpha = 0.35.
\label{eq:qstep}
\end{equation}

The complexity descriptor consists of the four decoder-reproducible
structural quantities above. Because all four structural quantities come
from anchor positions, entropy-model predictions, and the already-decoded
mask, the decoder reconstructs exactly the same steps without any
training images or training-time statistics. The bitstream records only
the formula version, base steps, and network parameters in a small
header and carries no per-anchor step array; the benefits of
content-adaptive quantization are thus fully retained inside a
self-contained bitstream.

During training the complexity network is optimized jointly with the
base entropy model: formula quantization is enabled after the 20,000-th
iteration and smoothed over the following 10,000 iterations, and
quantization uses a straight-through estimator to keep the pipeline
differentiable. The complexity network is a small $4\rightarrow
25\rightarrow 3$ network, where 25 is the hidden-layer width and the
three outputs are the adjustment logits for the Feature, Scaling, and
Offset fields; its impact on model size is negligible.

The two innovations act on the conditional probability model and the
quantization step respectively, and neither modifies HAC++'s anchor
generation, Gaussian expansion, or rendering pipeline. All conditioning
information needed by the decoder is carried by the bitstream itself,
which is the basis for the method's deployability.

\subsection{Bitstream Contract and Standalone Decoding Consistency}
\label{sec:bitstream}
Bitstream size follows the decode-required payload convention: only
files actually needed by a clean decoder are counted (anchor-position
stream, Feature/Scaling/Offset streams, masks, hash parameters, MLP
weights, context payload, and headers), excluding checkpoints and debug
data.

The encoder performs anchor-geometry coding, arithmetic coding of
attribute fields, and context writing; the decoder reconstructs the
model from the bitstream alone in a fresh process. Consistency
verification compares context, entropy parameters, formula inputs,
quantization steps, masks, and decoded symbols between the two sides;
bit-level consistency is declared only when the mismatch counts of
symbols, quantization steps, and integer CDFs are all zero.

\begin{table}[t]
\caption{Implementation and evaluation settings}
\centering
\begin{tabular}{@{}ll@{}}
\toprule
\textbf{Item} & \textbf{Setting} \\
\midrule
Training iterations & 90,000 \\
Anchor growth stopped after & 45,000 iterations \\
Rate--distortion weight & \(\lambda = 0.004\) \\
Anchor feature dimension & 50 \\
Offsets per anchor & 10 \\
Fine-grid spacing & 0.001 \\
Reported quality metrics & Decoded PSNR / SSIM / VGG-LPIPS \\
Size convention & Decode-required bitstream payload \\
\bottomrule
\end{tabular}
\label{tab:settings}
\end{table}

\subsection{Relationship to HAC/HAC++}
HAC \cite{chen2024} and HAC++ \cite{chen2025} predict anchor-attribute
distributions from hash-grid context and are the direct basis of this
work. We keep their Anchor--Hash main pipeline and mask-aware probability
modeling. The differences are two-fold. The context is extended from
single-scale fine features to a hierarchical structure of fine features,
parent features, and distance labels. Quantization steps are generated by
a decoder-reproducible formula rather than transmitted as additional
data. Neither module changes anchor generation, keyframe
structure, or the inference rendering path, so the proposal can be viewed
as two enhancements to HAC++ that directly reuse its open-source
implementation and deployment pipeline at low migration cost.

\subsection{Implementation and Complexity}
For large scenes, context construction runs in chunks to avoid memory
and kernel pressure; local-density estimation uses a deterministic
uniformly sampled subset when the anchor count is large, so both ends
stay reproducible. Beyond the anchor-attribute streams, the method adds
only 2 bits per anchor for the distance label plus a light header, so
the extra payload is negligible.

\section{Experimental Results and Analysis}
\subsection{Dataset and Experimental Setup}
We evaluate on three representative large-scale UAV scenes, 1-100, 3-07,
and 4-28, which differ in scale span, flight organization, and detail
density. All experiments train for 90,000 iterations with
rate--distortion weight $\lambda \in \{0.001, 0.002, 0.004, 0.006\}$.
Size is measured by decode-required payload and quality by decoded
rendering results.

Inference efficiency uses same-card serial profiling with 10 warm-up
frames followed by 150 test views over 3 rounds (450 timed frames).

\subsection{Rate--Distortion Results}
\begin{table*}[t]
\caption{Operating points under four rate--distortion weights
(PSNR / decode-required bitstream MB)}
\centering
\small
\begin{tabular}{@{}l cccc@{}}
\toprule
\textbf{Scene} & \(\lambda=0.001\) & \(\lambda=0.002\) &
\(\lambda=0.004\) & \(\lambda=0.006\) \\
\textbf{} & \textbf{PSNR / Size} & \textbf{PSNR / Size} &
\textbf{PSNR / Size} & \textbf{PSNR / Size} \\
\midrule
1-100 & 29.750 / 30.180 & 29.231 / 23.010 & 29.296 / 17.151 &
28.386 / 14.438 \\
3-07 & 28.333 / 27.850 & 27.342 / 21.347 & 28.629 / 15.488 &
26.730 / 12.400 \\
4-28 & 28.760 / 13.133 & 28.634 / 9.531 & 28.563 / 6.355 &
28.062 / 5.147 \\
\bottomrule
\end{tabular}
\label{tab:rd}
\end{table*}
Table~\ref{tab:rd} reports 12 rate--distortion operating points across
four $\lambda$ values. At the default $\lambda=0.004$, the
decode-required bitstreams are 17.151/15.488/6.355 MB with PSNR
29.296/28.629/28.563 dB for 1-100/3-07/4-28. Compared with the HAC++
reference points, 4-28 improves PSNR by 0.25 dB while reducing bitstream
size by 8.5\%, 1-100 attains its best PSNR (29.296 dB), and 3-07 reduces
bitstream size by 11.4\% at near-parity quality; at the more size-oriented
$\lambda=0.006$, the three scenes reduce bitstream size further by about
14.0\%/29.0\%/25.9\%. These gains are scene dependent and do not establish
a uniform rate--distortion advantage over HAC++.

Table~\ref{tab:rd} also shows that as $\lambda$ increases, bitstream size
decreases monotonically: for 1-100, size drops from 30.180 MB to 14.438
MB, while the default $\lambda=0.004$ point keeps a high PSNR of 29.296
dB at 17.151 MB. The method thus offers a range of operating points from
high fidelity to high compression, letting deployments choose
freely under bandwidth and quality budgets.

\subsection{Rate--Distortion Analysis}
\label{sec:rd}
We present the complete rate--distortion curve instead of a single
operating point, and quantify the comparison with HAC++ through
Bjontegaard metrics so that the reader can see the full behavior rather
than a scene-picked sample. Two qualitative findings hold already.

\begin{itemize}
    \item \textbf{The curve is available, not a singleton.} Across the
    sweep, higher $\lambda$ monotonically reduces bitstream size while
    degrading PSNR; the default point sits on this curve as a
    representative, not an isolated result. This turns the evaluation
    into a full quality--rate tradeoff rather than a ranked snapshot.
    \item \textbf{The comparison with HAC++ is reported, not assumed.}
    The per-scene outcomes are scene dependent: our method reaches its
    best PSNR on two scenes and a smaller size on the third, but these
    do not constitute a uniform rate--distortion advantage. We report
    the matched-budget BD-Rate below and do not claim global dominance.
\end{itemize}

\paragraph{BD-Rate / BD-PSNR (matched budget).}
We compute Bjontegaard delta on the matched-budget HAC++ and SHARC-GS
curves, with Pareto/monotone handling so that non-monotonic operating
points are integrated only after a valid frontier is formed. The
three-metric averages and per-scene values are:
\begin{center}\textbf{[TODO: BD-Rate / BD-PSNR table -- PSNR / SSIM / LPIPS, per-scene + macro]}
% TODO(B1): fill from compute_bdrate.py once the matched-budget curve is ready.
\end{center}

\paragraph{Equal-rate quality comparison.}
Because a single per-scene PSNR does not isolate rate efficiency, we
interpolate our curve at each HAC++ operating-point bitrate and compare
quality at equal size (and, conversely, size at equal quality).
\begin{center}\textbf{[TODO: equal-rate interpolation table -- Ours@HAC++.bpp vs HAC++]}
% TODO(m3): fill once RD curve is available.
\end{center}

\paragraph{Per-component bitstream breakdown.}
To make the size convention explicit, we report the decode-required
payload decomposed into anchor location, feature, scaling, offsets,
masks, hash parameters, MLP weights, and context.
\begin{center}\textbf{[TODO: per-component size table -- MB per component per scene]}
% TODO(A5): fill from codec_summary.json once the wave produces it.
\end{center}

\paragraph{Multi-seed stability.}
The default sweep is a single run. We therefore repeat the main
configuration with three seeds and report mean $\pm$ standard deviation,
so that any $\lambda$-ordering (and the reported non-monotonicity) can be
attributed to a stable effect rather than to training noise.
\begin{center}\textbf{[TODO: mean$\pm$std RD figure -- 3 seeds per operating point]}
% TODO(B2): fill after multi-seed wave.
\end{center}

\subsection{Ablation Study}
\label{sec:ablation}
To verify the independent contribution of each module, we run controlled
switch ablations on all three scenes: with $\lambda=0.004$ and the
90,000-iteration budget fixed, only the hierarchical-context and
formula-quantization switches vary, and the all-off configuration is the
baseline. Table~\ref{tab:ablation} summarizes the macro averages: I1
contributes mainly quality (+0.16 dB), I2 contributes mainly size
($-11.5\%$), and the joint configuration achieves the highest PSNR gain
in the table (+0.28 dB) while retaining most of the size reduction
($-8.5\%$). Because I1 changes the conditional distribution, I2
re-distributes its rate allocation in the joint run. This is an
operating-point choice that trades about 3\% of size for 0.05 dB of
quality relative to I2 alone. On the detail-rich 1-100 scene the
combination brings the largest gain (+0.76 dB, $-9.6\%$). A purely
size-oriented deployment can keep I2 only ($-11.5\%$). In every
formula-mode run the per-anchor step payload is zero, which shows that
reproducible quantization introduces no additional side information.

\begin{table}[t]
\caption{Ablation study (\(\lambda=0.004\), three-scene macro average,
relative to the all-off baseline)}
\centering
\begin{tabular}{@{}l cc@{}}
\toprule
\textbf{Configuration} & \(\Delta\)\textbf{PSNR (dB)} &
\(\Delta\)\textbf{Size (\%)} \\
\midrule
I1 only (hierarchical context) & +0.160 & $-1.2$ \\
I2 only (formula quantization) & +0.229 & $-11.5$ \\
I1 + I2 & +0.276 & $-8.5$ \\
\bottomrule
\end{tabular}
\label{tab:ablation}
\end{table}

\begin{figure}[t]
\centering
\includegraphics[width=0.92\linewidth]{fig2.png}
\caption{Ablation results (left: \(\Delta\)PSNR; right: size change;
three-scene macro average).}
\label{fig:ablation}
\end{figure}

\subsection{Quantitative Comparison with Existing Methods}
Table~\ref{tab:cmp} compares our method with nine representative existing
methods per scene, ranked following the convention of the prior
compression survey. On 4-28 it reaches its best PSNR (28.563 dB) with a
model 8.5\% smaller than HAC++; on 1-100 it reaches 29.296 dB at a size
comparable to HAC++; and on 3-07 it reaches near-HAC++ quality with an
11.4\% smaller model. These operating points are scene dependent and do
not establish a uniform rate--distortion advantage.

Across the three scenes, the two modules keep a favorable quality--size
trade-off---its best PSNR on two scenes and an 11.4\% smaller model than
HAC++ on the third, subject to the scene dependency noted above.

\begin{table*}[t]
\centering
\caption{\textbf{Quantitative comparison with existing methods.} We report decoded PSNR (dB, $\uparrow$), SSIM ($\uparrow$), perceptual LPIPS ($\downarrow$), and the decode-required bitstream size in MB ($\downarrow$) for SHARC-GS and nine representative methods on the 4-28, 1-100, and 3-07 scenes. Methods are listed in the ordering of the prior compression survey; bold, underlined, and yellow cells mark the best, second, and third value per metric per scene.}
\label{tab:cmp}
\resizebox{\linewidth}{!}{%
\begin{tabular}{l*{12}{c}}
\toprule
\multirow{2}{*}{\textbf{Method}}
& \multicolumn{4}{c}{\textbf{4-28}}
& \multicolumn{4}{c}{\textbf{1-100}}
& \multicolumn{4}{c}{\textbf{3-07}} \\
\cmidrule(lr){2-5}\cmidrule(lr){6-9}\cmidrule(lr){10-13}
& PSNR\(\uparrow\) & SSIM\(\uparrow\) & LPIPS\(\downarrow\) & Size\(\downarrow\)
& PSNR\(\uparrow\) & SSIM\(\uparrow\) & LPIPS\(\downarrow\) & Size\(\downarrow\)
& PSNR\(\uparrow\) & SSIM\(\uparrow\) & LPIPS\(\downarrow\) & Size\(\downarrow\) \\
\midrule
SHARC-GS (Ours)
& \best{28.563} & 0.888 & 0.298 & 6.355
& \best{29.296} & \second{0.905} & \second{0.137} & 17.151
& \second{28.629} & \second{0.839} & \second{0.222} & 15.488 \\
HAC++~\cite{chen2025}
& \second{28.311} & 0.890 & 0.293 & 6.946
& \second{28.998} & \best{0.912} & \best{0.118} & 16.787
& \best{28.772} & \best{0.860} & \best{0.187} & 17.473 \\
CAT
& 27.468 & 0.886 & \best{0.184} & \third{5.261}
& 26.291 & 0.830 & 0.175 & \second{13.114}
& 26.765 & 0.770 & 0.271 & \third{12.530} \\
PC-GS
& 27.324 & 0.886 & 0.300 & \second{5.040}
& 26.335 & 0.838 & 0.237 & \best{11.549}
& 26.971 & 0.788 & 0.296 & \second{12.161} \\
contextgs
& 27.602 & \second{0.894} & \third{0.289} & 8.652
& 27.752 & 0.871 & 0.196 & 30.165
& \third{28.104} & 0.828 & 0.249 & 23.656 \\
HAC~\cite{chen2024}
& 27.598 & \third{0.892} & 0.292 & 9.387
& 26.041 & 0.835 & 0.234 & 21.692
& 27.477 & 0.803 & 0.279 & 21.779 \\
RDO-gaussian
& 25.424 & 0.864 & 0.320 & \best{2.975}
& 26.099 & 0.824 & 0.253 & \third{13.383}
& 24.857 & 0.720 & 0.367 & \best{6.012} \\
TC-GS
& 27.304 & 0.885 & 0.302 & 7.327
& 26.431 & 0.826 & 0.261 & 18.429
& 26.117 & 0.748 & 0.343 & 16.243 \\
Scaffold-gs~\cite{lu2024}
& \third{27.677} & \best{0.898} & \second{0.281} & 132.520
& \third{28.394} & \third{0.889} & \third{0.166} & 287.151
& 27.648 & \third{0.836} & \third{0.235} & 247.447 \\
compact3d
& 25.872 & 0.879 & 0.307 & 8.510
& 26.528 & 0.841 & 0.242 & 25.880
& 25.673 & 0.754 & 0.341 & 16.420 \\
\bottomrule
\end{tabular}}
\end{table*}

\subsection{Qualitative Results}
Fig.~\ref{fig:qual} shows the qualitative comparison and per-pixel error
maps at the default operating point on the three scenes.

\begin{figure*}[t]
\centering
\includegraphics[width=0.92\textwidth]{fig3.png}
\caption{Qualitative comparison on the three scenes: GT, HAC++, Ours,
and per-method error maps against GT.}
\label{fig:qual}
\end{figure*}

\subsection{Standalone Decoding Consistency}
For the three default $\lambda=0.004$ bitstreams, only decode-required
files are copied to a clean directory and decoded in a fresh process. On
all three scenes the mismatch counts of symbols, quantization steps, and
integer CDFs are zero, which is consistent with a fully self-contained,
bit-exact bitstream; the distance label is a lossless 2-bit integer,
so the two sides share exactly the same conditioning, and its payload is
about 0.12/0.12/0.06 MB---far smaller than the attribute streams
themselves.

\subsection{Mechanism Analysis}
\label{sec:mechanism}
Beyond aggregate quality, we report evidence for \emph{why} the two
modules help, so that the gains can be attributed to a mechanism rather
than to a picked operating point.

\paragraph{Does content-adaptive quantization track spatial detail?}
We test whether the decoder-reproducible descriptor $u_i$ is larger on
high-frequency and reflective regions (building edges, fine texture,
water/specular highlights) than on flat regions. If this holds, the
formula quantizer is allocating finer steps where error is visually
weighted, rather than perturbing the step uniformly. We report
high-band PSNR (frequency/Laplacian decomposition), per-region PSNR,
and the descriptor distribution against region content.
\begin{center}\textbf{[TODO: descriptor-vs-high-frequency figure + per-region PSNR table]}
% TODO(R-Q2): fill once decoded renders are available.
\end{center}

\paragraph{Is the bitstream stable across views and time?}
A decoder-reproducible codec must not trade single-view PSNR for
multi-view or temporal stability. We render a smooth fly-through from
the test poses and measure temporal LPIPS, per-pixel flicker, optical-flow
warped error, and multi-view reprojection consistency for both SHARC-GS
and HAC++ under the same trajectory.
\begin{center}\textbf{[TODO: temporal/multi-view consistency table + figure]}
% TODO(M3): fill once the trajectory rendering pipeline runs.
\end{center}
If our method shows higher flicker in some regions, we report this as a
trade-off rather than claiming parity.

\paragraph{Do I1 and I2 interact positively?}
We report the multi-seed interaction of I2-only against I1+I2, so that
the joint run's small size penalty (about 3\% here) can be judged as a
stable operating-point choice rather than an unexplained loss.
\begin{center}\textbf{[TODO: I1/I2 interaction scatter + multi-seed error bars]}
% TODO(m2): fill after the multi-seed wave.
\end{center}

\subsection{Inference Efficiency}
Three-scene averages of same-card profiling: parameters decrease from
71.45M to 47.07M ($-34.12\%$), mean rendering latency from 16.48 to
14.43 ms ($-12.47\%$), P95 latency from 18.98 to 15.85 ms ($-16.48\%$),
peak allocated memory from 4168.56 to 3853.83 MB ($-7.55\%$), and frame
rate improves from 62.73 to 71.73 FPS ($+14.36\%$); encoding and
decoding times drop by 19.37\% and 21.76\% respectively. The compressed
model therefore improves inference efficiency; the results are
summarized in Table~\ref{tab:efficiency}.

\begin{table}[t]
\caption{Inference-efficiency comparison (same-card three-scene average)}
\centering
\footnotesize
\begin{tabular}{@{}l c c c@{}}
\toprule
\textbf{Metric} & \textbf{Ref.\ model} & \textbf{Ours} &
\textbf{Change} \\
\midrule
Parameters (M) & 71.45 & 47.07 & $-34.12\%$ \\
Mean render latency (ms) & 16.48 & 14.43 & $-12.47\%$ \\
P95 latency (ms) & 18.98 & 15.85 & $-16.48\%$ \\
Peak allocated memory (MB) & 4168.56 & 3853.83 & $-7.55\%$ \\
Frame rate (FPS) & 62.73 & 71.73 & $+14.36\%$ \\
Encoding time (s) & 39.43 & 31.80 & $-19.37\%$ \\
Decoding time (s) & 57.17 & 44.73 & $-21.76\%$ \\
\bottomrule
\end{tabular}
\label{tab:efficiency}
\end{table}

\subsection{Future Work}
Future work includes extending the hierarchical context to more levels
and larger scenes, exploring mask-aware progressive decoding orders for
streaming, and validating transferability on more public datasets. These
directions can all be realized incrementally without changing the
current bitstream main pipeline.

\section{Conclusion}
This contribution proposes SHARC-GS on the HAC++ framework: a
scale-aware hierarchical context lets the entropy model exploit
cross-scale structural priors, and decoder-reproducible formula
quantization brings content-adaptive quantization into practice without
side information. The two modules leave the HAC++ main pipeline
unchanged, all bitstreams are fully self-contained, and bit-exact
standalone decoding verification is passed. On three large-scale UAV
scenes the method reaches its best decoded PSNR on two scenes and a smaller
bitstream than HAC++ on the third, balancing fidelity, size, and
reproducibility, while the rate--distortion outcomes remain scene
dependent. The result is a step toward deployable compression for
large-scale UAV scenarios.

From a standardization perspective, the self-contained bitstream,
decoder-reproducible quantization rule, and bit-exact consistency
verification provide a clear interoperability contract across
implementations; the two modules are minimally invasive to the HAC++
pipeline and can be adopted incrementally on existing compression
frameworks. This contribution is a step toward practical compression of
3D representations of large-scale UAV scenes.

\section*{Acknowledgment}
The authors thank the anonymous reviewers and the supporting team for
their helpful comments.

\begin{thebibliography}{00}
\bibitem{kerbl2023} B. Kerbl, G. Kopanas, T. Leimk\"uhler, and G.
Drettakis, ``3D Gaussian splatting for real-time radiance field
rendering,'' \emph{ACM Transactions on Graphics}, vol. 42, no. 4, pp.
139:1--139:14, 2023.
\bibitem{chen2024} Y. Chen, Q. Wu, W. Lin, et al., ``HAC: Hash-grid
assisted context for 3D Gaussian splatting compression,'' in \emph{Proc.
European Conf. Computer Vision (ECCV)}, 2024.
\bibitem{chen2025} Y. Chen, Q. Wu, W. Lin, et al., ``HAC++: Towards 100X
compression of 3D Gaussian splatting,'' arXiv preprint
arXiv:2501.12255, 2025.
\bibitem{lu2024} T. Lu, H. Yu, L. Chen, et al., ``Scaffold-GS: Structured
3D Gaussians for view-adaptive rendering,'' in \emph{Proc. IEEE/CVF Conf.
Computer Vision and Pattern Recognition (CVPR)}, 2024, pp.
20654--20664.
\end{thebibliography}

\end{document}
